% ===================================================================
%  TAULA N° 3 — Enfança e Joventut.
%  Traduction 2026 d'apres texto/io, controlee sur texto/fr et
%  texto/ca. Consigne : voir texto/oc/10-taula-01.tex.
% ===================================================================

%%K t03-apar-1 apar
\VUsaut{3.0mm}
\VUtitre{15.0pt}{60}{TAULA N° 3}
\VUsaut{1.6mm}
\VUtitre{11.4pt}{0}{Enfança e Joventut.}
\VUsaut{3.4mm}

%%K t03-apar-2 apar
(Las taulas 3 e 4 son dispausadas de manièra que presentan, dins quatre
\VUgras{escènas familialas}, lo vocabulari essencial del \VUgras{temps}, de
l'\VUgras{atge}, de la \VUgras{familha}, de la \VUgras{rauba} e de la
\VUgras{santat}.)
\VUsaut{4.0mm}

%%K t03-tit-1 sub
\VUcentre{12.0pt}{0}{\textit{Primièra scèna.}}
\VUsaut{3.0mm}

%%K t03-tit-2 sub
\VUpk{10.2pt}{40}{Un batejal al vilatge.}
\VUsaut{2.4mm}

%%K t03-tit-3 sub
\VUpk{10.2pt}{40}{La prima.}
\VUsaut{3.0mm}

%%K t03-c1-01-1 p
1. --- Sèm a la \VUgras{prima}, al mes de \VUgras{març}, d'\VUgras{abril} o de
\VUgras{mai}, un jorn de la \VUgras{setmana} --- \VUgras{diluns},
\VUgras{dimars} o \VUgras{dimècres}. Es \VUgras{dètz} oras del matin.

%%K t03-c1-02-1 p
2. --- Fa bon \VUgras{temps}, l'\VUgras{aire} es pur. Lo solelh lusís. De
\VUgras{nivolets} \textsuperscript{(2)} pujan pel \VUgras{cèl}
\textsuperscript{(1)} blau.

%%K t03-c1-03-1 p
3. --- Los \VUgras{arbres} an tot bèl just de \VUgras{fuèlhas}. Los prims
\VUgras{pibols} \textsuperscript{(3)} e l'aut \VUgras{platan}
\textsuperscript{(17)} pòrtan pas encara que de borrons.

%%K t03-c1-04-1 p
4. --- Las \VUgras{dindoletas} \textsuperscript{(4)} son tornadas e volastrejan amb
de crits alègres a l'entorn de la \VUgras{flècha} \textsuperscript{(7)} del
\VUgras{cloquièr} \textsuperscript{(5)} que corona la \VUgras{glèisa}
\textsuperscript{(6)} del vilatge.

%%K t03-tit-4 sub
\VUsaut{3.4mm}
\VUpk{10.2pt}{40}{La Glèisa.}
\VUsaut{3.0mm}

%%K t03-c1-05-1 p
5. --- Fòrça simpla es aquesta glèisa amb son grand \VUgras{portal}
\textsuperscript{(13)} e son gròs \VUgras{relòtge} \textsuperscript{(8)}.
L'\VUgras{agulha pichona} \textsuperscript{(9)} es sus la chifra X: marca las
\VUgras{oras}. La \VUgras{granda} \textsuperscript{(10)} es sus las XII
(miègjorn); marca las \VUgras{minutas}.

%%K t03-c1-06-1 p
6. --- Dins lo cloquièr, las \VUgras{campanas} \textsuperscript{(11)} sonan
alègrament. Amont sul teulat s'auça una pichona \VUgras{crotz}
\textsuperscript{(12)} de pèira.

%%K t03-c1-07-1 p
7. --- La glèisa a pas solament un grand \VUgras{portal} \textsuperscript{(13)}, a
tanben una pòrta pichona de costat. Per aquesta pòrta lo sénher \VUgras{curat}
\textsuperscript{(15)}, vestit de sa \VUgras{sotana}, sortís de la
\VUgras{sagrestia} \textsuperscript{(14)} e se dirigís cap a la \VUgras{clastra}
\textsuperscript{(19)}.

%%K t03-c1-08-1 p
8. --- A prèp d'el, vaicí la \VUgras{lachièira} \textsuperscript{(24)} amb son
\VUgras{ase} \textsuperscript{(23)}. Torna de la vila, ont a portat son bon lach
pels enfants.

%%K t03-tit-5 sub
\VUsaut{4.0mm}
\VUpk{10.2pt}{40}{Lo vergièr.}
\VUsaut{3.4mm}

%%K t03-c1-09-1 p
9. --- Lo sénher Aliboron (l'ase) es fòrça content d'aquesta corsa matinala.
Aspira amb delici l'aire embaumat pel perfum de las \VUgras{flors}
\textsuperscript{(20)} que cobrisson los \VUgras{arbres fruchièrs}
\textsuperscript{(18)} del \VUgras{vergièr} vesin. Tot ròses e blancs son aqueles
\VUgras{persegièrs} \textsuperscript{(18)} e aqueles \VUgras{abricotièrs}
\textsuperscript{(19)}, que prometon una bona meissa.

%%K t03-c1-10-1 p
10. --- Pendent aqueste temps, las pichonas \VUgras{abelhas}
\textsuperscript{(22)} del \VUgras{bornhon} \textsuperscript{(21)} i van libar son
doç \VUgras{mèl}, en bronzinant.

%%K t03-c1-11-1 p
11. --- Mas qué se passa? Vaicí los enfants del vilatge, qu'arriban en corrent e
fan fugir las \VUgras{aucas} \textsuperscript{(25)} espaurugadas. Es lo seguici que
sortís de la glèisa aprèp lo \VUgras{batejal} del filh del notari.

%%K t03-c1-12-1 p
12. --- Lo pichon Jaume, dins son \VUgras{brèç} \textsuperscript{(26)}, se
desrevelha al son alègre de las campanas, e sa sòrre Magdalena, que vòl tanben
veire lo seguici, demanda a sa maire de venir la penchenar e la vestir sul lindal
de l'ostal.

%%K t03-c1-13-1 p
13. --- La maire i consentís. Se pausa son \VUgras{mocador de cap}
\textsuperscript{(28)}, son \VUgras{corset} \textsuperscript{(29)} e son
\VUgras{davantal} \textsuperscript{(30)}, e ven s'assetar sus una cadièireta
davant la pòrta. Magdalena pòrta pas que sos \VUgras{jupons}
\textsuperscript{(31)} e son \VUgras{cotelh} \textsuperscript{(32)}.

%%K t03-c1-14-1 p
14. --- Qu'es urosa! Oblida sas flors preferidas, sos \VUgras{clavèls}
\textsuperscript{(34)}, ont se pausan las \VUgras{parpalhòlas}
\textsuperscript{(35)}, sas \VUgras{tulipas} \textsuperscript{(36)}, sos
\VUgras{muguets} \textsuperscript{(37)} dins de \VUgras{tèsts}
\textsuperscript{(38)} sul \VUgras{mur} \textsuperscript{(33)}, e sas
\VUgras{ròsas} \textsuperscript{(39)}, que fan una tan polida bartassa. Agacha las
riquas raubas de las dònas del seguici.

%%K t03-c1-15-1 p
15. --- Los dròlles del vilatge fan pas gaire cas de las raubas. Se precipitan e
s'empenhon per aver de \VUgras{bonbons} \textsuperscript{(55)} o de
\VUgras{moneda} \textsuperscript{(52)}. Un d'eles plora \textsuperscript{(40)}.

%%K t03-c1-16-1 p
16. --- Un autre a caucat la pata del \VUgras{gos} \textsuperscript{(42)}, que
fugís en cridant e fa fugir la \VUgras{galina} \textsuperscript{(43)} e sos
\VUgras{polets} \textsuperscript{(44)}.

%%K t03-tit-6 sub
\VUsaut{5.0mm}
\VUpk{10.2pt}{40}{Lo seguici del batejal.}
\VUsaut{4.0mm}

%%K t03-c1-17-1 p
17. --- Davant lo seguici camina la \VUgras{noirissa} \textsuperscript{(45)}.
Pòrta una polida \VUgras{còfa} \textsuperscript{(46)} amb de rubans longs. Pòrta
dins sos braces lo \VUgras{nòvel-nascut} \textsuperscript{(47)}, tot vestit de
\VUgras{dentèlas} \textsuperscript{(48)}.

%%K t03-c1-18-1 p
18. --- Darrièr la noirissa venon lo \VUgras{pairin} \textsuperscript{(49)} e la
\VUgras{mairina} \textsuperscript{(53)}. Repartisson los bonbons e la moneda. Lo
pairin pòrta de \VUgras{gants} \textsuperscript{(51)} e un \VUgras{capèl de
tuba} \textsuperscript{(50)}. La mairina pòrta a la man un polit \VUgras{ramelet}
\textsuperscript{(54)}.

%%K t03-c1-19-1 p
19. --- Darrièr eles vesèm l'\VUgras{oncle} \textsuperscript{(56)} e la
\VUgras{tanta} \textsuperscript{(58)} de l'enfant. La tanta pòrta un elegant
\VUgras{capèl} \textsuperscript{(59)}, l'oncle una \VUgras{levita}
\textsuperscript{(57)} negra e un gilet blanc. Venon puèi lo \VUgras{cosin}
\textsuperscript{(60)} e la \VUgras{cosina} \textsuperscript{(61)}. Aquesta ten per
la man la \VUgras{sòrre} \textsuperscript{(62)} del nòvel-nascut, la pichona
Berta.

%%K t03-c1-20-1 p
20. --- L'autre \VUgras{fraire} \textsuperscript{(63)} de Berta camina darrièr.
Dona la man a una \VUgras{parenta} \textsuperscript{(64)}. Los \VUgras{amics}
\textsuperscript{(65)} de la familha venon aprèp.

%%K t03-tit-7 sub
\VUsaut{4.6mm}
\VUpk{10.2pt}{40}{Los que gaitan.}
\VUsaut{3.4mm}

%%K t03-c1-21-1 p
21. --- A prèp del seguici, vaicí un \VUgras{mendicant} \textsuperscript{(66)}
apiejat sus sa \VUgras{crossa} \textsuperscript{(67)}. Demanda l'almòina.

%%K t03-c1-22-1 p
22. --- De l'autre costat i a un dròlle fòrça \VUgras{cortés}
\textsuperscript{(68)}. Saluda e dona lo \VUgras{«bonjorn»} a las personas que
coneis.

%%K t03-c1-23-1 p
23. --- Los del vilatge son sortits de l'ostal per veire passar lo seguici del
batejal. Vesèm un \VUgras{païsan} \textsuperscript{(69)} amb sa \VUgras{bonèta de
coton} e a son costat una \VUgras{païsana} \textsuperscript{(70)}. Puèi una
\VUgras{vièlha} \textsuperscript{(71)} apiejada sus son \VUgras{baston}
\textsuperscript{(72)}. Enfin, lo \VUgras{molinièr} \textsuperscript{(73)} amb son
larg \VUgras{capèl de feutre} \textsuperscript{(74)}, e una jove païsana caussada
d'\VUgras{esclòps} \textsuperscript{(75)}, qu'ensenha de caminar a son enfant.

%%K t03-c1-24-1 p
24. --- Es una scèna fòrça divertenta.
\VUsaut{4.0mm}

%%K t03-tit-8 sub
\VUcentre{12.0pt}{0}{\textit{Segonda scèna.}}
\VUsaut{3.0mm}

%%K t03-tit-9 sub
\VUpk{10.2pt}{40}{Una fèsta publica.}
\VUsaut{2.4mm}

%%K t03-tit-10 sub
\VUpk{10.2pt}{40}{Jòcs nautics e regatas.}
\VUsaut{3.0mm}

%%K t03-c2-01-1 p
1. --- L'\VUgras{estiu} es arribat. Lo solelh ardent del mes de \VUgras{junh}
(julhet o \VUgras{agost}) convida los joves a se banhar e a sortir en barca.

%%K t03-c2-02-1 p
2. --- Sus la riba drecha del riu, los remaires se despolhan per prene part als
jòcs nautics e a las regatas.

%%K t03-c2-03-1 p
3. --- Un d'eles se lèva las \VUgras{sabatas} \textsuperscript{(1)}; las
desliga (las desbotona). Sos dos companhs son ja prèstes. Pòrtan pas que son
\VUgras{jersei} \textsuperscript{(2)} e son \VUgras{calçon de banh}
\textsuperscript{(3)}. Charran en esperant sos amics. Parlan de la fièira e de la
fèsta que se tenon sus l'autra riba.

%%K t03-c2-04-1 p
4. --- Per tèrra, a lor costat, jai la \VUgras{rauba} de sos companhs: un
\VUgras{parelh} de \VUgras{botinas} \textsuperscript{(4)}, de \VUgras{sabatas}
\textsuperscript{(5)}, de \VUgras{caucetas} \textsuperscript{(6)}, de capèls de
\VUgras{feutre} \textsuperscript{(7)} e de \VUgras{palha} \textsuperscript{(8)} e
una \VUgras{cravata} \textsuperscript{(9)}.

%%K t03-c2-05-1 p
5. --- Un dels joves a plegat amb suènh sa \VUgras{vèsta} \textsuperscript{(10)}.
La pausa sus una toalha, dins la quala son vesin embolharà lèu los dos
\VUgras{vestits}.

%%K t03-c2-06-1 p
6. --- Un seisen dròlle es en retard. Pòrta encara sa \VUgras{casqueta}
\textsuperscript{(11)} sul cap. S'es levat ni la \VUgras{camisa}
\textsuperscript{(12)}, ni lo \VUgras{còl} \textsuperscript{(13)}, ni los
\VUgras{ponhets} \textsuperscript{(14)}. Ten a la man son \VUgras{gilet}
\textsuperscript{(17)} e pòrta encara sas \VUgras{bragas} \textsuperscript{(16)} e
sas \VUgras{bretèlas} \textsuperscript{(15)}. Segur que serà pas prèst per la
corsa que ven.

%%K t03-c2-07-1 p
7. --- A prèp dels joves, un caminòl bordat de \VUgras{cardons}
\textsuperscript{(18)} davala a la riba del riu, ont l'oncle Jaume prepara, entre
las \VUgras{canas} \textsuperscript{(19)}, los \VUgras{fuòcs artificials}
\textsuperscript{(20)} que se tiraràn al vèspre.

%%K t03-tit-11 sub
\VUsaut{4.4mm}
\VUpk{10.2pt}{40}{La corsa.}
\VUsaut{3.4mm}

%%K t03-c2-08-1 p
8. --- Sul riu, qualques dònas que s'abrigan jos sos parasolelhs se passejan en
\VUgras{barca} \textsuperscript{(21)}. Seguisson la corsa, qu'es fòrça
interessanta. Dins cada \VUgras{iòla} \textsuperscript{(22)}, quatre
\VUgras{remaires} \textsuperscript{(24)} raman de totas lors fòrças, mentre que lo
\VUgras{timonièr} \textsuperscript{(26)} ten lo \VUgras{timon}
\textsuperscript{(25)} e dirigís la frala embarcacion. Los \VUgras{rems}
\textsuperscript{(23)} tustan l'aiga a la mesura, e las barcas leugièras navègan
a una velocitat verginosa.

%%K t03-c2-09-1 p
9. --- Qualques joves se disputan, contra lo pilar del pont, lo prètz de la
\VUgras{cocanha orizontala} \textsuperscript{(28)}. Un d'eles avança amb precaucion
sul pal sople e lisant per atendre la pichona \VUgras{bandièra}
\textsuperscript{(29)} fixada a la cima. Son malastruc \VUgras{concurrent}
\textsuperscript{(30)} es tombat dins l'aiga. Nada per arribar a la riba.

%%K t03-c2-10-1 p
10. --- De dròlles mai grands \VUgras{jostan en barca} \textsuperscript{(32)}
contra lo \VUgras{cai} \textsuperscript{(33)}. A totes aqueles jòcs assistís un
\VUgras{public} \textsuperscript{(31)} nombrós, apiejat sul \VUgras{parapet} del
\VUgras{pont} \textsuperscript{(27)} e sarrat sus la \VUgras{riba}. Qualques
espectators, çaquelà, se viran per seguir lo jòc de la \VUgras{cocanha}
\textsuperscript{(45)} o per admirar lo \VUgras{seguici municipal} que ven a la
\VUgras{distribucion} dels prètzes.

%%K t03-c2-11-1 p
11. --- D'en primièr, vaicí de \VUgras{piòts de carrièra} \textsuperscript{(35)}
que van davant los \VUgras{musicaires} \textsuperscript{(36)}; puèi venon lo
\VUgras{cònsol} \textsuperscript{(37)}, sos adjonchs e lo \VUgras{conselh
municipal} de la pichona vila. Tampan la marcha los \VUgras{pompièrs}
\textsuperscript{(38)}.

%%K t03-tit-12 sub
\VUsaut{5.0mm}
\VUpk{10.2pt}{40}{La fièira.}
\VUsaut{3.6mm}

%%K t03-c2-12-1 p
12. --- I a fòrça mond sul \VUgras{camp de fièira} \textsuperscript{(39)}. Se ven
veire las diferentas \VUgras{barracas}: los \VUgras{cavalets de fusta}
\textsuperscript{(40)}, lo \VUgras{circ} \textsuperscript{(41)}, ont la
representacion es ja començada; lo \VUgras{bal public} \textsuperscript{(42)}, ont
los joves e las jovas de la vila gaudisson del plaser de \VUgras{dançar}.

%%K t03-c2-13-1 p
13. --- Al fons vesèm las \VUgras{arenas} \textsuperscript{(44)}, ont lo valent
\VUgras{matador} espanhòl luta contra lo \VUgras{taur} \textsuperscript{(45)}
furiós; puèi la \VUgras{montanha russa} \textsuperscript{(49)} e la \VUgras{barraca
de tir} \textsuperscript{(46)}, ont s'exercisson a manejar las \VUgras{armas de
fuòc} e a tirar a la \VUgras{cibla}.

%%K t03-c2-14-1 p
14. --- A prèp d'aquí, vaicí lo \VUgras{teatre de fièira} \textsuperscript{(47)},
ont se representan la \VUgras{comèdia} e lo \VUgras{drama}. Fòrça prèp i a la
barraca del \VUgras{gigant} \textsuperscript{(48)} e del \VUgras{nan}. La
\VUgras{talha} del primièr es de dos mètres quinze; lo segond es tot bèl just tan
grand coma un enfant.

%%K t03-c2-15-1 p
15. --- Enfin, vaicí la granda \VUgras{menagariá} \textsuperscript{(50)} amb sos
\VUgras{animals sauvatges}, son \VUgras{tigre} \textsuperscript{(51)} de
poderosas \VUgras{grifas}, son \VUgras{leon} \textsuperscript{(52)} d'espessa
\VUgras{cabeladura}, sas doas \VUgras{girafas} \textsuperscript{(53)} e son
\VUgras{elefant} \textsuperscript{(54)}, que se sèrv amb tanta adreça de sa
\VUgras{trompa}.

%%K t03-c2-16-1 p
16. --- Lo \VUgras{domdaire} \textsuperscript{(55)} que dòmda aqueles animals es a
la pòrta de la barraca.
\VUsaut{6.0mm}
\VUfilet{22mm}
